\title{Byte Latent Transformer: Patches Scale Better Than Tokens}

\begin{abstract}
We present the Byte Latent Transformer ({\textsc BLT}), a novel byte-level large language model (LLM) architecture that achieves parity with tokenization-based LLMs in performance at scale, while substantially enhancing both inference efficiency and robustness. {\textsc BLT} operates by encoding input bytes into adaptively sized patches that become the central units for model computation. These patches are delineated by evaluating the entropy of the subsequent byte, thereby assigning greater model resources and computational effort to data segments exhibiting higher complexity. We conduct the first computationally-controlled (flop-controlled) scaling analysis of byte-level models, extending up to 8B parameters and 4T training bytes. Our findings illustrate that models trained directly on byte streams, eschewing a predetermined vocabulary, can indeed be scaled effectively. Efficiency in both training and inference is bolstered by the model's adaptive patch selection—opting for longer patches in segments characterized by high predictability—while also yielding qualitative benefits in reasoning and improved generalization on rare cases. Overall, at equivalent inference cost, {\textsc BLT} delivers superior scaling compared to models using fixed tokenization, achieved by jointly increasing both the size of the patches and the underlying model.
\end{abstract}

\section{Introduction}
\label{section:intro}

We present the Byte Latent Transformer~(\textbf{BLT}), an architecture that eliminates the need for tokenization by directly learning from unprocessed byte sequences. Remarkably, BLT achieves parity with tokenization-dependent models at scale, while delivering notable enhancements in both efficiency and robustness (\S\ref{sec:robustness}). In contrast, most state-of-the-art large language models (LLMs) operate in an almost fully end-to-end fashion, with the exception of the tokenization step—a rule-based preprocessing phase that clusters bytes into a predetermined vocabulary of tokens. This process introduces biases into string compression and is associated with limitations such as modality and domain-specific sensitivity~\citep{dagangetting}, heightened vulnerability to input noise~(\S\ref{sec:robustness}), insufficient handling of orthographic variations~\citep{edman2024cute}, and inequitable multilingual processing~\citep{liang2023xlm,petrov2024language,limisiewicz2024myte}.

Historically, tokenization has been indispensable because training llms directly on byte sequences entails prohibitive computational expense at scale, mainly owing to the increased sequence lengths~\citep{xue2022byt5}. Earlier research attempts to address this challenge have relied on more computationally efficient self-attention mechanisms~\citep{el2020characterbert, clark2022canine} or by adopting architectures that dispense with attention entirely~\citep{wang2024mambabyte}~(\S\ref{section:related_work}). Nonetheless, such strategies predominantly benefit the training of \textit{small models}. When scaling up, the primary computational burden of Transformers arises from the extensive feed-forward network layers applied to each byte, rather than from attention-related operations.

To optimize the distribution of computation, we introduce a flexible, trainable approach for aggregating bytes into \textit{patches}~(\S\ref{section:patching}), alongside a novel model architecture that integrates both byte-level and patch-level data. In contrast to tokenization, {\textsc BLT} does not utilize a predetermined set of patch types. Instead, any combination of bytes can be transformed into latent patch embeddings by means of lightweight, learned encoder and decoder components. Our findings indicate that this strategy yields a \textit{greater} efficiency in compute allocation compared to models that rely on tokenization.

Tokenization-based LLMs devote equal computational resources to every token, regardless of individual complexity. This approach sacrifices efficiency for overall performance because tokenization relies on compression heuristics, which do not necessarily align with the actual difficulty of making predictions. A core concept of our proposed architecture is that computation should be adaptively allocated according to need. For instance, predicting the endings of most words generally constitutes an easy, low-entropy task, and thus does not require the full capacity of a large transformer, in contrast to generating the initial word of a sentence, which involves more uncertainty. This principle underpins the design of {\textsc BLT}'s architecture~(\S\ref{section:architecture}), which comprises three transformer components: a pair of lightweight byte-level \textit{local models} and a larger, global \textit{latent transformer}~(\autoref{fig:arch_high_level}).

In order to decide how bytes should be aggregated into patches—and thus determine a dynamic distribution of computational effort—{\textsc BLT} utilizes the entropy of the next-byte prediction for segmentation. This results in groups of bytes that share similar contextual information density, supporting a more balanced and efficient computational process.

This work introduces the first scaling analysis of byte-level models constrained by floating-point operation counts, covering architectures with up to 8B parameters trained on 4T bytes, and demonstrates that it is feasible to train such models from raw bytes at full scale without relying on a predetermined token vocabulary. Overall, {\textsc BLT} achieves training performance on par with Llama 3, as measured under controlled computational budgets,\footnote{Computational costs are measured by counting the total number of floating-point operations (flops) performed by the model.} but requires as much as 50\% fewer flops during inference (\S\ref{section:scaling}).
Moreover, we find that operating directly on byte sequences leads to notable gains in capturing long-tail data distributions. {\textsc BLT} models exhibit greater resilience to noisy textual inputs when compared to their tokenizer-dependent counterparts and show improved capacity for character-level processing, which is substantiated through tasks targeting orthographic variation, phonological structure, and machine translation for under-resourced languages (\S\ref{sec:robustness}).
Additionally, the {\textsc BLT} architecture allows for joint scaling of both the model and input patch sizes without increasing the inference compute budget. Expanding the patch size can lower per-inference computational requirements on average, which permits allocating saved compute toward expanding the latent transformer—since its execution frequency decreases. Our flop-controlled scaling experiments at inference time (\autoref{fig:fixed_inference_scaling}) demonstrate distinctly improved scaling properties in comparison to approaches reliant on traditional tokenization.

To conclude, this work offers several key contributions: 1) We present {\textsc BLT}, a byte-level latent LLM architecture that enables dynamic compute allocation to boost flop efficiency; 2) We demonstrate that our approach can match Llama 3 in terms of training flop budget up to the 8B parameter scale, with the added flexibility to exchange small decreases in evaluation metrics for up to 50\% improvements in flop efficiency; 3) {\textsc BLT} facilitates a novel paradigm in scaling LLMs, permitting increases in model size without exceeding a predetermined inference computation budget; 4) We provide evidence of {\textsc BLT} models’ superior robustness to input noise and enhanced sensitivity to sub-word properties of the input, aspects which are often overlooked by traditional token-based LLMs. The training and inference code for {\textsc BLT} is made available at~\url{https://github.com/facebookresearch/blt}.

\section{Patching: From Individual Bytes to Groups of Bytes}
\label{section:patching}

Dividing byte sequences into \textit{patches} enables {\textsc BLT} to flexibly allocate computational resources according to contextual requirements. Figure~\ref{fig:patching_types} illustrates a variety of approaches for partitioning byte sequences into patches. In formal terms, a patching function $f_p$ partitions an input byte sequence $\pmb{x} = \{x_i, |i=1, \ldots, n\}$ of length $n$ into $m < n$ patches, denoted $\pmb{p} = \{p_j|j=1,\ldots,m\}$, by associating each $x_i$ with an indicator in the set \{0, 1\}, where a value of 1 signals the initiation of a new patch. For both token-level and patch-level architectures, overall computation is largely dictated by the number of processing steps taken by the principal Transformer block. In the {\textsc BLT} framework, this step count corresponds to the total patches produced by the patching function when encoding input data. Thus, the average quantity of bytes per patch—referred to as \textit{patch size}—serves as the central determinant for computational effort during training and inference under any specific patching scheme~(\S\ref{section:flops}).

We proceed by outlining three specific patching strategies: segmenting input into patches of uniform byte count~(\S\ref{section:static-patch}), utilizing whitespace as patch boundaries~(\S\ref{section:white-patch}), and dynamically constructing patches by leveraging entropy estimates from a lightweight byte language model~(\S\ref{section:dyn-patch}). Subsequently, we elaborate on the concept of incremental patching and clarify distinctions between patching and standard tokenization approaches~(\S\ref{section:bpe-patch}).

\subsection{Strided Patching Every K Bytes}
\label{section:static-patch}

A simple approach to organizing bytes is to divide them into uniformly sized patches of length $k$, as demonstrated in MegaByte~\citep{yu2023megabyte}. This method employs a fixed stride, making both implementation and scaling during training or inference straightforward, as well as allowing for uncomplicated adjustments to the average patch size, and thus provides an easy way to manage computational cost. Nevertheless, this type of patching introduces notable limitations. One issue is the lack of adaptability in computation allocation: processing steps may be wasted on uninformative segments (such as spaces in code) while insufficient compute may be dedicated to segments containing highly concentrated information (e.g., mathematical expressions). Additionally, this strategy can result in inconsistent and non-context-aware segmentations of similar byte sequences—for example, causing the same word to be divided in varying ways across contexts.

\subsection{Space Patching}
\label{section:white-patch}

According to \citet{slagle2024spacebyte}, a straightforward yet powerful refinement to strided patching is introduced whereby new patches are formed immediately following any space-like byte\footnote{
    Space-like bytes refer to those bytes that are neither Latin letters, digits, nor utf-8 continuation bytes. Furthermore, every patch must include at least one non space-like byte.}, taking advantage of the fact that such bytes often demarcate linguistic boundaries across many languages. In Space patching, an additional transformer inference step (involving increased computational flops) is devoted to every word. This approach standardizes the patching of words across different sequences and directs computational effort toward challenging predictions, which frequently occur right after spaces. For instance, identifying the initial byte in the response to the question ``Who composed the Magic Flute?\uline{\hspace{.6em}}'' presents a much greater prediction challenge than filling in subsequent bytes after the first ``M'', since determining the first character substantially narrows the set of probable answers, rendering the rest of ``Mozart'' relatively easier to guess. Nonetheless, space patching exhibits limitations: it does not efficiently accommodate all languages and domains and, critically, it lacks the capacity to adjust patch size. To address this, we proceed to present a novel patching approach that capitalizes on the insight that the initial bytes of words tend to be the most unpredictable, while also furnishing an intuitive method for dynamically adapting patch size.

\subsection{Entropy Patching: Using Next-Byte Entropies from a Small Byte LM}
\label{section:dyn-patch}

Instead of utilizing a heuristic approach like identifying whitespace, we adopt a data-driven methodology to pinpoint instances where the next-byte prediction is most uncertain. To this end, we propose \textit{entropy patching}, a technique that leverages entropy estimates to determine patch boundaries.

We train a small byte-level auto-regressive language model on the training data for {\textsc BLT} and compute next byte entropies under the LM distribution $p_{e}$ over the byte vocabulary $\mathcal{V}$:

\begin{equation}
H(x_i) = \sum_{v \in \mathcal{V}} p_{e}(x_i=v|\pmb{x}_{<i}) \log p_{e}(x_i=v|\pmb{x}_{<i})
\end{equation}

We employ two strategies for determining patch boundaries using entropies $H(x_i)$. The first approach selects points where the entropy exceeds a fixed global threshold, as depicted in~\autoref{fig:patching}. The second method detects points whose entropy is notably higher than that of the preceding point. This latter technique can alternatively be viewed as locating positions where the nearly monotonic decline in entropy within a patch is interrupted.

\begin{align*}
    \text{Global Constraint}\hspace{1cm}H(x_t)& > \theta_g\\
    \text{Approx. Monotonic Constraint}\hspace{1cm}H(x_t)& - H(x_{t-1}) > \theta_r
\end{align*}

Patch boundaries are determined in an efficient preprocessing phase that takes place while loading data. This contrasts with the approach of~\citet{nawrot-etal-2023-efficient}, where a classifier is employed to detect entropy-driven patch boundaries. Within our experiments~(\S\ref{section:experiments}), we evaluate and compare both strategies for separating low-entropy and high-entropy bytes.

\subsection{The Byte-Pair Encoding (BPE) Tokenizer and Incremental Patching}
\label{section:incremental-patching}
\label{section:bpe-patch}

Numerous contemporary llms, such as our baseline Llama 3, rely on subword tokenizers, for example, bpe~\citep{gage1994new, sennrich-etal-2016-neural}. Throughout this work, we use the term ``tokens'' to denote byte groups selected from a \textit{finite} vocabulary established before training begins, distinguishing them from ``patches,'' which represent sequences formed dynamically, independent of a predefined vocabulary. An essential distinction between tokens and patches is that, when utilizing tokens, the underlying byte-level features are not directly accessible to the model.

A significant advancement introduced by {\textsc BLT} compared to tokenization-based models is its reconceptualization of the balance between vocabulary size and computational demands. In conventional llms, expanding the vocabulary results in a greater average token length, reducing the number of prediction steps required, yet simultaneously enlarging the output space of the final projection layer. This trade-off restricts the capacity of tokenization-dependent methods to produce meaningful shifts in both token granularity and the efficiency of inference. As an illustration, Llama 3 increases the mean token length from 3.7 to 4.4 bytes, but does so at the expense of quadrupling its embedding matrix relative to Llama 2.

During generation, {\textsc BLT} must determine at each byte position whether it has reached a patch boundary, since this dictates whether further processing is required using the Latent Transformer. Importantly, this decision must be made without knowledge of subsequent bytes that have not yet been produced; therefore, the patching procedure cannot rely on future elements of the sequence to guide segmentation. More precisely, a patching strategy $f_p$ possesses the characteristic of incremental patching if the following condition holds:
$$f_p(\pmb{x}_{<i}) = f_p(\pmb{x})_{<i}$$
It should be noted that bpe does not conform to the requirements of an incremental patching scheme because an identical initial subsequence may be segmented differently depending on what tokens follow, thus failing to meet the stated criterion\footnote{Introducing a dedicated boundary token to mark patch breaks enables bpe to serve as an incremental patching method, albeit with the tradeoff of lengthening the byte sequence.}.

\section{{\textsc BLT} Architecture}
\label{section:architecture}
{\textsc BLT} consists of a substantial global autoregressive language model functioning over patch-level representations, as well as a pair of more compact local models responsible for transforming byte sequences into patch representations and subsequently reconstructing bytes from those patches~(\autoref{fig:arch_high_level}).

\subsection{Latent Global Transformer Model}

\textit{The Latent Global Transformer} refers to an autoregressive transformer architecture $\mathcal{G}$ with $l_{\mathcal{G}}$ layers, designed to transform an input sequence of latent patch embeddings, $p_j$, into corresponding output patch embeddings, $o_j$. In this work, we employ the subscript $j$ to indicate patch indices and $i$ for byte positions. The global transformer adopts a block-causal attention scheme~\citep{dubey2024llama}, ensuring that each patch can attend only to itself and preceding patches within the same document. This global component is responsible for most of the floating point operations both in pre-training and inference stages. Therefore, being selective about its application enables us to adjust and tailor computational resource allocation across various segments of the input and output, in accordance with their respective complexities.

\subsection{Local Encoder}
\label{section:local}

\textit{The Local Encoder Model}, represented as $\mathcal{E}$, consists of a compact transformer architecture, featuring $l_{\mathcal{E}} << l_{\mathcal{G}}$ layers, and is primarily responsible for rapidly converting a sequence of input bytes $b_i$ into rich patch-level representations, $p_j$. Unlike traditional transformers, this model integrates a cross-attention layer following each transformer block, enabling the aggregation of byte-level features into patch-level embeddings~(\autoref{fig:crossattn}).
Initially, the byte sequence $b_i$ is mapped via an embedding matrix of shape $\mathbb{R}^{256 \times h_{\mathcal{E}}}$, yielding the embedded vectors $x_i$. These representations may further incorporate supplementary data by means of hash-embeddings~(\S\ref{sec:hash-emb}).
Subsequently, a stack of transformer and cross-attention layers~(\S\ref{sec:enc-cross-attn}) is applied to successively refine and condense these embeddings into patch representations, $p_i$, which are then forwarded to the global transformer, $\mathcal{G}$, for additional processing.
Within this architecture, transformer layers operate under a \textit{local block causal} attention paradigm: each token is permitted to attend only to the $w_{\mathcal{E}}$ tokens immediately preceding it within a context window. Although this window may overlap dynamic patch divisions, it does not extend beyond the limits of the document. Subsequent subsections elaborate on the specifics of the embedding process and the mechanics of the cross-attention component.

\subsubsection{Encoder Hash n-gram Embeddings}
\label{sec:ngram-emb}
\label{sec:hash-emb}
An essential aspect of building strong and nuanced representations at each timestep $i$ is embedding knowledge of the prior sequence of bytes. In {\textsc BLT}, this is accomplished by simultaneously modeling the current byte $b_i$ on its own, as well as incorporating its membership within a byte n-gram context. Specifically, at every position $i$, we initially form byte-grams

\begin{align}
    g_{i,n}=\{b_{i-n + 1},\ldots, b_i\}
\end{align}

for each byte position $i$ and for values of $n$ ranging from three to eight.\footnote{
We disregard $n$-gram features of length $n$ or greater whenever $i<n$. }

Subsequently, we propose hash $n$-gram embeddings, where every byte $n$-gram is projected through a hash function onto a specific slot of a fixed-size embedding matrix $E_{n}^{hash}$ for each $n \in\{3,4,5,6,7,8\}$~\citep{bai2010learning}. The embedding obtained from this process is combined with the corresponding byte embedding, after which normalization occurs and the representation is forwarded to the local encoder architecture. The computation of the enriched embedding is given by

\begin{align}
e_i &= x_i + \sum_{n = 3,...,8}E_{n}^{hash}(\text{Hash}(g_{i,n})) \\
\text{where, Hash}(g_{i,n}) &= \text{RollPolyHash}(g_{i,n}) \% |E_{n}^{hash}|
\end{align}

We adjust $e_i$ by dividing it by the total count of $n$-gram sizes incremented by one, employing RollPolyHash as specified in Appendix~\ref{appendix:rollpolyhash}. Section~\ref{section:ablations} presents an ablation analysis on the impact of $n$-gram hash embeddings under varying $n$ values and embedding table dimensions, particularly examining trends under flop constraints. Alongside the hashed $n$-gram embeddings, our investigations also encompassed frequency-based $n$-gram embeddings; further information on these studies is provided in Appendix~\ref{appendix:ngrams}.

\subsubsection{Encoder Multi-Headed Cross-Attention}
\label{sec:enc-cross-attn}
Our approach is largely based on the input cross-attention mechanism introduced in the Perceiver architecture~\citep{jaegle2021perceiver}. However, our method diverges in that latent variables are dynamically associated with patch representations, rather than being derived from a predefined set of latent vectors~(\autoref{fig:crossattn}), and each latent representation is restricted to attending only to the bytes present in its corresponding patch. For each patch $p_j$, we construct a query vector by first aggregating the byte-level representations for that patch through a pooling operation. This pooled vector is then transformed via a linear projection, specifically $\mathcal{E}_{C} \in \mathbb{R}^{h_{\mathcal{E}} \times (h_{\mathcal{E}} \times U_{\mathcal{E}})}$, where $U_{\mathcal{E}}$ denotes the number of heads in the encoder’s cross-attention. Formally, defining $f_{\text{bytes}}(p_j)$ as the sequence of bytes linked to patch $p_j$, the computation proceeds as follows:

\begin{align}
P_{0,j} &= \mathcal{E}_{C}(f_\text{bytes}((p_j)), f~ \text{is a pooling function} \\
P_l &= P_{l-1} + W_o\left(\text{softmax}\left(\frac{QK^T}{\sqrt{d_k}}\right)V\right) \\
\text{where } Q_j &= W_q(P_{l-1,j}), K_i = W_k(h_{l-1,i}), V_i = W_v(h_{l-1,i}) \\
h_l &= \text{Encoder-Transformer-Layer}_l(h_{l-1})
\end{align}

Here, $P \in \mathbb{R}^{n_p \times h_{\mathcal{G}}}$ denotes the matrix containing representations of $n_p$ patches that are forwarded to the global model. The initialization of these patch vectors is achieved by aggregating the byte embeddings $e_i$ associated with each respective patch $p_j$. The matrices $W_q$, $W_k$, $W_v$, and $W_o$ perform the linear projections for queries, keys, values, and output, where both the keys and values are derived by projecting the byte representations $h_i$ from the preceding layer ($e_i$ is used in the base layer). Our approach involves a targeted masking technique that ensures each query $Q_j$ focuses exclusively on those keys and values connected to the bytes comprising patch $j$. Since multi-headed attention is utilized across $Q$, $K$, and $V$—with patch representations often possessing higher dimensionality ($h_{\mathcal{G}}$) relative to $h_{\mathcal{E}}$—we structure $P_l$ as multiple attention heads of size $h_{\mathcal{E}}$ in the cross-attention operation, and later concatenate these head outputs to reach the final $h_{\mathcal{G}}$ dimensionality. Furthermore, a pre-LayerNorm is applied individually to the queries, keys, and values. Note that this cross-attention layer omits positional embeddings. Finally, the output of the cross-attention module is integrated with its input via a residual connection.

\subsection{Local Decoder} 

Analogous to the local encoder, the local decoder $\mathcal{D}$ is constructed as a lightweight transformer architecture, consisting of $l_{\mathcal{D}}$ layers, where $l_{\mathcal{D}} << l_{\mathcal{G}}$. This component translates a sequence of global patch representations $o_j$ into raw byte outputs $y_i$. The process of predicting each byte depends on the previously decoded bytes, with the decoder utilizing the hidden states generated by the local encoder corresponding to the byte sequence. Across its $l_{\mathcal{D}}$ layers, the architecture alternates between cross-attention and transformer operations. Within each layer, cross-attention precedes the transformer sublayer, enabling the construction of byte representations from the input patch representations; the subsequent local decoder transformer layer then processes this resulting sequence of bytes.

\subsubsection{Decoder Multi-headed Cross-Attention} 

In the decoder cross-attention, the roles of the queries and key/values are interchanged i.e. the byte-representations are now the queries, and the patch representations are now the key/values. The initial byte-representations for the cross-attention are initialized as the byte embeddings from the last encoder layer i.e. $h_{l_{\mathcal{E}}}$. The subsequent byte-representations for layer $l$, $d_{l,i}$ are computed as:

\begin{align}
D_0 &= h_{l_{\mathcal{E}}} \\
B_l &= D_{l-1} + W_o\left(\text{softmax}\left(\frac{QK^T}{\sqrt{d_k}}\right)V\right), \\
  &\text{where } Q_i = W_q(d_{l-1,i}), K_i = W_k(\mathcal{D}_C(o_j)), V_i = W_v(\mathcal{D}_C(o_j)) \\
  D_l &= \text{Decoder-Transformer-layer}_l(B_l)
\end{align}

Here, as before, $W_k$ and $W_v$ denote the key and value projection matrices, each performing a linear transformation followed by the split operation $\mathcal{D}_C$ on the global model’s final patch representations $o_j$. The matrix $W_q$ is used to project the query, operating on byte representations $d_{l-1}$ from the previous layer of the decoder transformer (or on $h_{l_{\mathcal{E}}}$ in the initial layer). $W_o$ refers to the output projection matrix, so that $B$ has dimensions $\mathbb{R}^{h_{\mathcal{D}} \times n_b}$, with $n_b$ representing the number of output bytes. To obtain the subsequent decoder representations $D_l$, a decoder transformer layer processes the output $B$ from the cross-attention step. Analogous to the cross-attention in the local encoder, we employ multi-headed attention, apply pre-LayerNorm, forgo positional embeddings, and implement a residual connection surrounding the cross-attention unit.

\section{Experimental Setup}
\label{section:experiments}
We conduct systematic experiments to directly evaluate {\textsc BLT} against models based on tokenization, making deliberate efforts to avoid conferring any undue advantage to {\textsc BLT} due to potentially longer input context lengths.

\subsection{Pre-training Datasets} In this work, models at all evaluated scales are initially trained on two primary datasets: (1) The Llama 2 dataset~\citep{touvron2023llama}, containing 2 trillion tokens sourced from a wide range of publicly available datasets that have undergone extensive cleaning and filtering processes to enhance quality; and (2) {\textsc BLT}-1T: a novel dataset comprising 1 trillion tokens, aggregated from multiple public repositories and also incorporating a selection of the Datacomp-LM pre-training data~\citep{li2024datacomp}. The Llama 2 dataset supports scaling law analyses—particularly to establish the optimal token count, as discussed by~\cite{dubey2024llama}—which guide architectural decisions for {\textsc BLT}. In contrast, the {\textsc BLT}-1T dataset is employed for full pre-training to enable performance comparisons against Llama 3 on various downstream evaluations. It is noteworthy that neither dataset contains data extracted from Meta platforms or services. For baseline assessments utilizing tokenizer-based approaches, we adopt the Llama 3 tokenizer, configured with a 128K token vocabulary. This tokenizer, in our empirical tests, yielded consistently superior baseline outcomes compared to the tokenizer from Llama 2.

\subsection{Entropy Model} For our experiments, the entropy model utilized is a byte-level language model that is trained on an identical data distribution to the full {\textsc BLT} model. Unless specified otherwise, the architecture comprises a transformer with 100 million parameters, organized in 14 layers with a hidden size of 512, and it applies sliding window attention over 512 bytes. All other hyperparameter settings follow those used in our local and global transformer configurations. We further assessed variants with different model capacities, context ranges, and model designs, as outlined in \autoref{section:ablations}. Notably, when the context window of the model is sufficiently narrow, the resulting entropy model can be compactly represented using an efficient lookup table.

\subsection{Entropy Threshold and Equalizing Context Length} For approaches that utilize entropy-based patching, we determine a threshold for patching such that a targeted average \textit{patch size} is attained over the pretraining data mix. In {\textsc BLT}, as opposed to traditional tokenization, the \textit{patch size} is flexibly set, significantly affecting the amount of context provided to the model. To guarantee comparable average context length and to prevent models with larger patch sizes from gaining an undue advantage, we fix the expected total byte count in each batch. Consequently, for models with greater patch sizes, the sequence length is proportionally decreased. On Llama 2 data, an 8k byte context is utilized, whereas for the {\textsc BLT}-1T dataset, the context is increased to an average of 16k bytes, all while holding the average batch size steady at 16M bytes.

Although the average batch size remains fixed, dynamic patching techniques produce varying proportions of bytes per patch when assembling data batches. In our {\textsc BLT} training setup, patches within batches are aggregated so as to bypass unnecessary padding operations in the resource-intensive latent transformer; as a result, each batch contains an equal number of patches. To prevent memory surges caused by sequences with abnormally large patches, we pad—or when necessary, truncate—input byte sequences to lengths of 12k for Llama 2 and 24k for {\textsc BLT}-1T datasets throughout training.

\subsection{Entropy Model Context}
\label{section:entropy-context}
Through observation, we note that the application of entropy patching results in increasingly large patches, particularly within structured datasets such as multiple choice tasks (as illustrated by the patching process on an MMLU sample in \autoref{fig:mmlu_patching}). These datasets typically exhibit considerable repetition. The larger patch sizes are attributed to diminished entropy associated with redundant segments in the entropy model context. Therefore, in the main experiment with {\textsc BLT}-Entropy (patch size 4.5), we refresh the entropy context by inserting new lines and employ an approximate monotonicity constraint; this approach is less susceptible to "entropy drift" arising from shifts in context length. Notably, this adjustment pertains solely to the computation of entropy values—the overall methodology for determining the entropy threshold remains unchanged.

\subsection{FLOPs Estimation} 
\label{section:flops}

Our approach primarily adheres to the methodology outlined in Chinchilla~\citep{hoffmann2022training} for calculating transformer FLOPs, which includes operations from the feed-forward layers, qkvo projections within the self-attention mechanism, and both attention score computation and output projections. One significant distinction is our assumption that the input embedding layer operates as a highly efficient lookup table rather than performing dense matrix multiplication, thus incurring no FLOPs. Consistent with prior research, we approximate the FLOP count for the backward pass to be double that of the forward computation.

To compute flops \textit{per byte} for {\textsc BLT} models, we add up the flops for the local encoder transformer, the global latent transformer, and the local decoder transformer, together with the cross attention blocks in the encoder and the decoder:

\begin{align}
    \text{FL}_{\text{{\textsc BLT}}} &= \text{Transf. FL}(h_{\mathcal{G}}, l_{\mathcal{G}}, m=n_{ctx}/n_p, V=0) / n_p\\
    &\quad + \text{Transf. FL}(h_{\mathcal{E}}, l_{\mathcal{E}}, m=w_{\mathcal{E}}, V=0) \\
    &\quad + \text{Transf. FL}(h_{\mathcal{D}}, l_{\mathcal{D}}, m=w_{\mathcal{D}}, V=256) \\
    &\quad + \text{Cross Attn. FL}(h_{\mathcal{E}}, l_{\mathcal{E}}, m=n_p, r=n_p/k) \cdot k/n_p \\
    &\quad + \text{Cross Attn. FL}(h_{\mathcal{D}}, l_{\mathcal{D}}, m=k, r=k/n_p)
\end{align}

Here, $n_{ctx}$ denotes the sequence length measured in bytes, while $n_p$ indicates the patch size. The parameter $r$ represents the ratio between the number of queries and key/value pairs. The variable $k$ is defined as the quotient of the patch dimension by the byte dimension—this signifies the count of local model partitions that, when concatenated, make up the full global model embedding (in \autoref{fig:crossattn}, $k=2$). $V$ designates the vocabulary size that pertains solely to the output projection within the local decoder. The attention mechanism's effective context length, $m$, varies depending on whether it operates over the byte-level sequence or the patch-level sequence. We adjust the computation of attention flops for each component to account for these differences. Detailed formulas for the calculation of Transformer-FLOPs and Cross-Attention FLOPs are included in Appendix~\ref{section:flops-appendix}.

\subsection{Bits-Per-Byte Estimation} Perplexity only makes sense in the context of a fixed tokenizer as it is a measure of the uncertainty for each token. When comparing byte and token-level models, following previous work~\citep{xue2022byt5, yu2023megabyte, wang2024mambabyte}, we instead report Bits-Per-Byte (BPB), a tokenizer independent version of perplexity. Specifically:

\begin{align}
    \text{BPB}(x)&= \frac{\mathcal{L}_{CE}(\pmb{x})}{\ln(2)\cdot n_{\text{bytes}}}
\end{align}

Here, the aggregate cross-entropy loss, which quantifies the uncertainty of the data sequence $\pmb{x}$, is divided by both the total byte count in $\pmb{x}$ and a normalization constant.

\subsection{Transformer Architecture Hyperparameters} All transformer layers within {\textsc BLT}, encompassing both the local and global components, are designed with reference to the Llama~3~\citep{dubey2024llama} blueprint. The feed-forward networks leverage the SwiGLU activation function~\citep{swiglu}, while rotary positional embeddings (RoPE)~\citep{RoPe}—configured with $\theta=500000$ as recommended by~\citep{xiong2024effective}—are included solely in self-attention layers. RMSNorm \citep{RMSNorm} is employed for normalizing layers. For self-attention operations with static attention patterns such as \textit{block causal} or \textit{fixed-window block causal}, we rely on Flash attention~\citep{dao2022flashattention}, adopting a window size of 512 for masks with fixed width. Flex Attention\footnote{\url{https://pytorch.org/blog/flexattention}} is utilized to efficiently support the cross-attention layers due to their use of dynamically generated, patch-dependent masks, enabling fused computation and substantial training acceleration.

\subsection{{\textsc BLT}-Specific Hyperparameters}

In investigating the performance of {\textsc BLT} models, our experimental analysis addresses two key aspects: scaling behavior and downstream task performance. Models with varying parameter counts—specifically 400M, 1B, 2B, 4B, and 8B—were included in these analyses. Details on the model architectural configurations are provided in Appendix~Table~\ref{tab:arch}.

Initialization of queries for the initial cross-attention block within the local encoder is performed using max-pooling. Across all {\textsc BLT} models, we utilize a hashing setup comprising $500,000$ hashes, a single hash function, and n-gram sizes selected from the interval 3 to 8. All models are trained with a consistent learning rate of $4e-4$, a value determined to be optimal after a hyperparameter search in the $1e-3$ to $1e-4$ range for both the 400M and 1B models, where token and {\textsc BLT} models both converged on the same learning rate.

For the analysis of scaling trends using Llama-2 data, recommended training batch sizes from~\cite{dubey2024llama}—or their byte-wise equivalents—are employed. The AdamW optimizer~\citep{adamw} is used during training with $\beta_1 = 0.9$, $\beta_2 = 0.95$, and an $\epsilon = 10^{-8}$. The learning rate is scheduled with a linear warm-up phase over 2000 steps, followed by a cosine decay to zero. Additionally, a weight decay factor of 0.1 is applied, and global gradient norms are clipped to a maximum of 1.0.

\section{Scaling Trends}
\label{section:scaling}

We provide a comprehensive overview of the scaling behavior of byte-level models, offering insights that can guide future scaling of {\textsc BLT} models. Our study is designed to overcome key gaps observed in earlier examinations of byte-level models by: (a) analyzing scaling dynamics within the context of compute-optimal training protocols; (b) training aligned 8B parameter models using substantial datasets (reaching up to 1T tokens or 4T bytes), and assessing their performance across various downstream applications; and (c) evaluating scaling properties under controlled inference cost conditions. In subsequent sections, we delve into particular benefits that arise from representing data as byte-sequences.

\subsection{Parameter Matched Compute Optimal Scaling Trends}
\label{section:parameter-matched-scaling}
We employ the Llama 2 dataset to train multiple \textit{compute-optimal} models, including both bpe and {\textsc BLT} variants, at four distinct scales ranging from 1B to 8B parameters. Language modeling performance is then evaluated and plotted as a function of training flops, using a representative subset of the training data mixture. In training the bpe models, we utilize the optimal parameter-to-token ratio, as identified in Llama~3~\citep{dubey2024llama}. This \textit{compute-optimal} regime is theoretically justified to maximize training set performance under a fixed compute budget~\citep{hoffmann2022training}, and thus serves as a strong baseline for comparative assessment. For every bpe model, we additionally train a paired {\textsc BLT} model on identical data, equipping each with a Latent Transformer that mirrors the architecture and scale of the bpe Transformer counterpart.

As shown in~\autoref{fig:scaling-compute-opt} (right), the performance of {\textsc BLT} models is consistently on par with or superior to that of bpe models, a pattern observed across increasing model sizes and flops. To our knowledge, {\textsc BLT} represents the first byte-level Transformer architecture to demonstrate scaling behaviors comparable to BPE-based models within compute-optimal conditions. These findings substantiate our hypothesis that the parameter-to-training-compute ratio optimal for bpe architectures also holds for {\textsc BLT}, or at the very least, does not deviate substantially.

Both enhancements to model architecture and the implementation of dynamic patching are essential for aligning with bpe scaling trajectories. In ~\autoref{fig:scaling-compute-opt} (left), we present a comparison between models utilizing space patching and Llama 3. For SpaceByte \citep{slagle2024spacebyte}, we estimate its performance using {\textsc BLT} configured with space-patching, excluding both n-gram embeddings and cross-attention mechanisms. While SpaceByte demonstrates a gain relative to Megabyte, there is still a considerable gap compared to Llama 3. \autoref{fig:scaling-compute-opt} (right) depicts the collective impact of architectural innovation and dynamic patching. At scale, {\textsc BLT} models achieve results competitive with leading tokenizer-based architectures such as Llama 3.

Additionally, we note that the selection of tokenizer significantly impacts the performance of tokenizer-based models; specifically, models utilizing the Llama-3 tokenizer achieve superior results compared to those employing the Llama-2 tokenizer when both are trained on identical datasets.

In summary, the {\textsc BLT} model positions itself between Llama 2 and Llama 3 in terms of architectural trends, especially when much larger patch sizes are employed. Whereas the bpe tokenizers utilized by Llama 2 and Llama 3 yield mean token sizes of 3.7 and 4.4 bytes respectively, {\textsc BLT} demonstrates comparable scaling behavior with average patch sizes of 6 and even up to 8 bytes. Since inference flops are inversely related to average patch size, selecting an 8-byte patch size translates into almost 50\% reduction in inference computation. Notably, models configured with larger patch sizes not only benefit in terms of efficiency but also tend to exhibit improved performance as both model and dataset sizes are scaled up. For example, {\textsc BLT} at a patch size of 8 bytes initially underperforms relative to bpe Llama 2 at the 1B parameter mark but surpasses it by the time the scale reaches 7B parameters. These findings indicate that larger patch sizes might yield even greater advantages at increased scales, hinting at the viability of experimenting with still larger patch sizes as model and compute resources continue to expand.

\subsection{Beyond Compute Optimal Task Evaluations}
\label{section:evals}
To further investigate scaling behaviors, we train an 8B {\textsc BLT} model that surpasses the compute optimal regime using the {\textsc BLT}-1T dataset, which is both larger and of higher quality. The model's capabilities are then assessed across a broad set of standard classification and generation benchmarks. For evaluation purposes, we focus on tasks encompassing common sense reasoning, general world knowledge, and code generation:
\paragraph{Classification tasks} are drawn from ARC-Easy (0-shot)~\citep{Eval_ARC}, Arc-Challenge (0-shot)~\citep{Eval_ARC}, HellaSwag (0-shot)~\citep{Eval_hellaswag}, PIQA (0-shot)~\citep{Eval_piqa}, and MMLU (5-shot)~\citep{hendrycks2020measuring}.  Evaluation is performed using a prompt-scoring approach, which computes the probability of each answer by evaluating the likelihood assigned to its corresponding choice token. Results are presented as average accuracy.
\paragraph{Coding related generation tasks:}  To gauge the model’s proficiency in Python code synthesis, we report pass@1 rates for both MBPP (3-shot)~\citep{Eval_MBPP} and HumanEval (0-shot)~\citep{Eval_HumanEval} benchmarks.

In \autoref{tab:evals}, we present a comparison among three models trained on the {\textsc BLT}-1T dataset: a model using the bpe Llama 3 tokenizer,\footnote{The Llama 3 tokenizer, featuring a 128k vocabulary, was selected because it outperforms the 32k vocabulary of Llama 2.} and two {\textsc BLT} model variants. The first variant adopts a space-patching approach ({\textsc BLT}-Space), while the second uses an entropy-based patching strategy ({\textsc BLT}-Entropy), incorporates an approximate monotonicity constraint, and resets the entropy model's context with new lines as outlined in~\autoref{section:entropy-context}. All models are trained using the same amount of computation (flops). For {\textsc BLT}-Entropy, we further adjust the entropy threshold during inference from 0.6 to 0.1, a modification that yields improved task performance, though it requires additional inference steps.

The {\textsc BLT}-Entropy model achieves superior results compared to the Llama 3 model on four of the seven evaluated tasks, despite both models being trained with an equivalent data volume in terms of bytes. This enhanced performance likely stems from two main factors: (1) more efficient utilization of computational resources during training, facilitated by dynamic patching, and (2) the explicit modeling of information at the byte level rather than at the token level.

In contrast, {\textsc BLT}-Space lags behind the Llama 3 tokenizer across all tasks except one, although its use of a larger mean patch size—6 bytes—substantially lowers inference flops. For reference, both the bpe model and the entropy-patching-based model maintain a similar mean patch size of about 4.5 bytes on the training data mixture. Given an identical training budget, the model with the larger patch size processes 30\% more data than the other two approaches, which could potentially cause {\textsc BLT} to deviate further from achieving compute-optimality.

\subsection{Patches Scale Better Than Tokens} In {\textsc BLT} models, it is possible to expand both the model parameters and the size of patches without altering the overall training or inference flops, nor requiring an increase in the training data. This flexibility to extend patch size distinguishes patch-based approaches from fixed-vocabulary token-based models, which are constrained by efficiency trade-offs (see Section~\ref{section:bpe-patch}). By employing larger patch sizes, computational demands are reduced, enabling more resources to be devoted toward enlarging the global latent transformer, as it needs to execute fewer times.

To evaluate whether increasing model size while reducing the number of steps on larger patches leads to improved performance compared to employing smaller models for more steps, we undertake a fixed inference scaling study. Utilizing models of 400 million and 3.6 billion parameters initialized with the Llama 2 tokenizer, we identify models with comparable FLOP counts that use the Llama 3 tokenizer, as well as {\textsc BLT}-Entropy models that process average patch sizes of 6 and 8 bytes in the training data mixture (refer to~\autoref{tab:fixed-inf-params} for specifics on the models). For models employing a patch size of 8, we opt for 3 encoder layers in place of 1. All models are trained under a range of training FLOP constraints.

\autoref{fig:fixed_inference_scaling} illustrates that {\textsc BLT} models demonstrate superior scaling behavior relative to tokenization-based models across both classes of inference flops. Although BPE-based architectures exhibit stronger performance when training resources are limited, they are soon overtaken by {\textsc BLT} models, once training moves past the compute-optimal threshold. In real-world settings, allocating additional resources to pretraining—despite it being a one-off expenditure—can yield models with markedly better performance under a fixed inference budget. This trend is exemplified by the 8B model class, such as Llama 3.1, which has undergone pretraining on data volumes up to two orders of magnitude greater than the compute-optimal quantity for its parameter count.

The point at which {\textsc BLT} surpasses token-based models has shifted marginally nearer to the compute-optimal threshold in the context of larger flop class models, decreasing from 3x to 2.5x the compute-optimal allocation. The model employing patch size 8 demonstrates a sharper scaling trajectory in this larger flop regime, surpassing alternative models more rapidly. As outlined in Section~\ref{section:parameter-matched-scaling}, models utilizing larger patch sizes increasingly approach the performance of BPE-based models as model scale increases. This phenomenon can be partly attributed to the diminishing proportion of total computation required by the byte-level Encoder and Decoder components, which scale less rapidly than the Latent Transformer. Notably, scaling the total number of parameters by a factor of 20—expanding from 400M to 8B—results in only an approximate doubling of {\textsc BLT}'s parameters dedicated to local modules. This distinction is crucial since expanding patch size impacts the computational budget exclusively within the patch Latent Transformer, leaving byte-level module flops largely unaffected. Accordingly, this explains why the model size ratio for {\textsc BLT}-Entropy with ps=8 rose from 1.6x to 1.7x relative to Llama 2 when moving toward larger-scale models.

In conclusion, our investigation into patch-length scaling reveals that the {\textsc BLT} patch-based architecture exhibits superior scaling behavior when both patch size and model size are jointly increased. Notably, this favorable scaling trend continues—and may even strengthen—as model capacity grows.

\section{Byte Modeling Improves Robustness} We further evaluate the robustness of the {\textsc BLT} architecture in comparison to token-based models that do not utilize byte-level signals, and introduce a method for transforming pretrained token models to operate on bytes.
\label{sec:robustness}

\subsection{Character-Level Tasks}

One of the initial incentives for developing byte-level models stemmed from their ability to handle noise at the byte level within inputs, as well as their intrinsic understanding of the internal makeup of tokens—a capability that conventional tokenizer-based models often lack. To assess these attributes, we conduct supplementary tests using benchmarks specifically designed to gauge resilience to input perturbations and to evaluate comprehension of character-level elements, encompassing both English and multiple languages, as well as numerical digits and phonemic symbols. The outcomes of these experiments are summarized in Table~\ref{tab:char_tasks}.

\paragraph{Noisy Data} To assess and compare the resilience of {\textsc BLT} and tokenizer-based approaches, we construct noisy variants of the benchmark classification datasets detailed in Section~\ref{section:evals}. Specifically, we use five separate character-level noise augmentation techniques: (a) \textit{AntSpeak}: Transforms the entire input into sequences of uppercase characters with spaces between each character. (b) \textit{Drop}: Randomly deletes 10\% of the characters from the text. (c) \textit{RandomCase}: Randomly switches the case of characters, resulting in half uppercase and half lowercase characters. (d) \textit{Repeat}: Selects 20\% of the characters at random and repeats them up to four times. (e) \textit{UpperCase}: Converts every character in the string to uppercase. Each noise method is applied separately either to the prompt, the completion, or both, forming distinct evaluation scenarios. Scores are then averaged across these conditions. Table \ref{tab:char_tasks} displays the outcomes on the noised HellaSwag~\citep{Eval_hellaswag} dataset, revealing that {\textsc BLT} consistently achieves greater robustness compared to tokenizer-based baselines, outperforming a similarly trained model by an average of 8 points, and even surpassing the Llama 3.1 model despite its substantially larger training corpus.

\paragraph{Phonology - Grapheme-to-Phoneme (G2P)} We evaluate {\textsc BLT}'s proficiency in converting a sequence of graphemes—letters denoting a word—into its corresponding phonemic transcription, reflecting pronunciation. Table \ref{tab:char_tasks} summarizes the performance of the G2P task under a 5-shot learning scenario, employing the Phonology Bench~\citep{suvarna2024phonologybench} dataset. The results indicate that {\textsc BLT} surpasses the baseline Llama 3 1T tokenizer-based model for this particular task.

\paragraph{CUTE}
To measure character-level comprehension, we conduct evaluations of {\textsc BLT} using the CUTE benchmark~\citep{edman2024cute}, which consists of multiple tasks grouped into three main types: compositional understanding, orthographic similarity, and sequence manipulation capabilities. This benchmark presents a substantial difficulty for most models that rely on tokenization, as these models tend to recognize token spellings but find it challenging to leverage this information for effective text manipulation. As indicated in Table~\ref{tab:char_tasks}, {\textsc BLT}-Entropy achieves a performance lead of over 25 points compared to the BPE-based Llama 3 models on this benchmark. Notably, our model excels at tasks requiring manipulation at the character level, reaching a 99.9\% accuracy rate in both spelling-related evaluations. These considerable gains, achieved despite {\textsc BLT} being trained on only one-sixteenth the data used for Llama 3.1, suggest that BPE-based architectures struggle to acquire fine-grained character information. Figure \ref{fig:cute_examples} showcases examples highlighting cases in which the Llama 3 tokenizer-based model fails but {\textsc BLT} delivers strong results. Word deletion and insertion are the only task categories where the BPE approach surpasses ours. Byte-level models may find word-level manipulations more complex, yet the difference in performance remains modest; constructing representations from characters upwards to words may, in fact, be more tractable than decomposing from words to characters. All evaluations use identical setups across tasks and original prompts from Huggingface. It should be noted that BPE models might improve with targeted prompt modifications.

\paragraph{Low Resource Machine Translation} We assess the performance of {\textsc BLT} on translation tasks involving both directions for twenty-one lower-resourced languages, each spanning six widely-used language families and diverse scripts, utilizing the FLORES-101 benchmark~\citep{goyal2022flores}. SentencePiece BLEU scores are presented in Table~\ref{tab:flores}. The findings show that {\textsc BLT} surpasses a system employing the Llama 3 tokenizer, offering a 2-point improvement for translations into English and a 0.5-point gain for translations from English. On well-established language pairs, {\textsc BLT} matches or marginally exceeds the performance of Llama 3. Crucially, for numerous language pairs within lesser-resourced families, {\textsc BLT} delivers superior results compared to Llama 3, highlighting the advantages of byte-level modeling in handling and generalizing to rare or complex byte sequences.

\subsection{Training {\textsc BLT} Using Llama 3 Initialization}
\label{sec:distillation}
We investigate a methodology whereby {\textsc BLT} models benefit from established pre-trained models that utilize tokenizers, thereby enhancing training efficiency and convergence speed. This is accomplished by setting the global transformer weights in the {\textsc BLT} framework to match those from a pre-trained Llama 3.1 model at the start of training. Following this initialization, we optimize the global transformer parameters at a learning rate that is one-tenth of that used for the local encoder and decoder networks, adhering to the recommended number of training steps for Llama 3. Performance metrics, juxtaposed with a baseline {\textsc BLT} in Table~\ref{tab:distill}, indicate a clear advantage: the {\textsc BLT} model initialized from Llama 3.1 outperforms both the original Llama 3 and standard {\textsc BLT} baselines, all trained under identical FLOPs constraints. Further, in comparison with our {\textsc BLT}-Entropy variant (detailed in Table~\ref{tab:evals}), which was trained on a much more extensive dataset comprising 1T tokens, the Llama 3.1-initialized {\textsc BLT} still surpasses it on the MMLU benchmark. This underscores the method's potential for dramatically lowering the computational requirements of model training.

This approach can be interpreted as converting models that rely on tokenizers into ones that operate without them, thereby transforming a pre-trained LLaMA 3.1 model into a {\textsc BLT} model. For a thorough assessment, we compare the original LLaMA 3.1, trained with 15T tokens (see Table \ref{tab:distill}), to its {\textsc BLT} counterpart. The model exhibits minor losses in performance on MMLU and HumanEval, but experiences notably greater degradation on additional tasks. These findings highlight the necessity for future research to better exploit the underlying pre-trained model, with a particular focus on optimizing data composition and other key hyperparameters.

\section{Ablations and Discussion}
\label{section:ablations}
Here, we analyze various ablation experiments that support our design decisions regarding the architecture of {\textsc BLT} as well as the selection of the patching methodology and hyperparameter values for the {\textsc BLT} model with 8B parameters, trained using the {\textsc BLT}-1T dataset.

\paragraph{Entropy Model Hyper-parameters} In order to analyze how the size of the entropy model and the length of the context window influence scaling behavior, we experiment with byte-level entropy transformer architectures spanning from 1 million to 100 million parameters, while systematically varying the context window length between 64 and 512 tokens. Using our $400m$ and $1b$ {\textsc BLT} models trained on the Llama-2 dataset, we generate training flop scaling law plots of bits-per-byte (bpb), as depicted in \autoref{fig:entropy_ablation}. The results demonstrate that increasing either the model's parameter count or its context window generally enhances scaling performance, although the improvements plateau once the model exceeds 50 million parameters.

\paragraph{Types of Patching}

We investigate the impact of the four patching methods outlined in Section~\ref{section:patching}, namely: (1) Strided Patching with strides set to 4 and 6, (2) Patching at whitespace boundaries, (3) BPE Tokenizer patching utilizing the Llama 3 tokenizer, and (4) Entropy-based patching leveraging a compact byte-level LLM.

Although dynamic patching shortens the actual sequence length, we normalize the sequence length across different patching approaches to ensure that each model encounters a comparable context window. During both training and evaluation, every model processes an equivalent expected byte count per sequence, thus eliminating possible confounds linked to modeling extended contexts. Results of these ablation studies are presented in Figure~\ref{fig:scaling-compute-opt}. The data demonstrate that all non-static patching methods provide superior performance compared to static patching, with space patching performing nearly as well as dynamic entropy-based patching.

In \autoref{tab:evals_ablation}, we report benchmark results for {\textsc BLT} models, contrasting the performance of tokenizer-based models, space patching, and entropy-based patching. All models are trained on the Llama 2 dataset for the number of steps determined to be optimal~\citep{dubey2024llama}. While space patching offers a more straightforward approach by foregoing the use of an entropy model during training, our analysis reveals that the improvements linked to entropy-based patching in scaling experiments~(Section~\ref{section:scaling}) are also reflected in downstream benchmark tasks.\footnote{Space patching results shown are from earlier runs lacking cross-attention, though similar results hold when cross-attention is included.}

\paragraph{Cross-Attention} 
As shown in~\autoref{tab:crossattn_ablations_1b}, we perform ablation experiments by introducing cross-attention at different layers within both the encoder and decoder components of {\textsc BLT}. For encoder cross-attention, we explore three strategies for initializing the queries: (1) using a shared learned embedding across all global states, (2) employing a hash-based embedding derived from the bytes within the patch, and (3) applying pooling to the hidden representations of the patch bytes generated by the encoder at the specific layer under consideration.

Our results indicate that integrating cross-attention into the \textit{decoder} yields the greatest performance benefits. Incorporating cross-attention within the encoder offers marginal gains, and these are observed primarily when employing pooling-based query initialization. Moreover, our analysis reveals that cross-attention provides notable advantages on the Common-Crawl dataset, with the effect being more pronounced as the patch size increases.

\paragraph{n-gram Hash Embeddings} 

We examine the impact of varying n-gram hash embedding vocabulary sizes—specifically, 0, 100K, 200K, and 400K—and report our findings in Table~\ref{tab:ngram_ablations_1b}. Our analysis reveals that incorporating hash embeddings consistently benefits all tested domains, with notably greater improvements observed for Wikipedia and Github (showing a 0.04 bpb improvement compared to 0.01 bpb at the 8B scale after 15k training steps). Increasing the number of hashes from 500K to 300K at the 8B scale only marginally impacted performance, with a difference of roughly 0.001 bpb after 15k steps. This suggests that hash embeddings are essential for elevating the performance of {\textsc BLT} models to be competitive with those using traditional tokenizers; nevertheless, beyond 300K hash entries, additional gains become minimal. Furthermore, our results indicate that the benefits provided by hash embeddings are largely additive with those from cross-attention, as they enhance model performance on distinct datasets.

\paragraph{Local Model Hyperparameters} Table \ref{tab:local_layers} presents an ablation study of different configurations for the number of layers in both the local encoder and decoder. The results indicate that, in combination with hash n-gram embeddings, {\textsc BLT} achieves strong performance when utilizing a minimal encoder—specifically, a single layer—while employing a more complex, deeper decoder.

\section{Related Work}
\label{section:related_work}

\textbf{Character-Level RNNs:} The task of Character Language Modeling has remained prominent since the advent of neural architectures~\citep{sutskever2011generating,mikolov2012subword,graves2013generating}, due to its innate capability to seamlessly handle out-of-vocabulary terms without necessitating back-off strategies. \cite{kim2016character} introduced a model architecture where only the input sequence is represented at the character level, utilizing convolutional and highway network layers that serve as input to LSTM-based RNNs. Their approach achieved parity with contemporary RNN-based state-of-the-art language models on English while exceeding their effectiveness on languages with rich morphological structures—highlighting an important merit of character-level language models. In another study, \cite{kenter2018byte} approached machine comprehension with byte-level LSTMs, which demonstrated superior performance compared to word-level architectures for Turkish and Russian, languages marked by complex morphology. Similarly, \cite{zhang2015character} employed convolutional models based on character representations for text classification tasks, surpassing word-based models in selected scenarios. \citet{chung2019hierarchical} developed hierarchical LSTM frameworks that apply boundary detectors at successive levels to automatically uncover underlying textual hierarchy, further refining the results in character-level language modeling. Additionally, ByteNet, proposed by \cite{kalchbrenner2016neural}, leveraged convolutional neural networks on character data in lieu of attention mechanisms for machine translation purposes.

\textbf{Character-Level Transformers:} The advent of transformer architectures leveraging attention mechanisms~\citep{vaswani2017attention} in conjunction with subword tokenization approaches~\citep{sennrich-etal-2016-neural} has substantially advanced the capabilities of neural models on both language modeling tasks and standard benchmarks. Nevertheless, the use of word and subword tokens inherently imposes an inductive bias regarding the abstraction level at which such models function. To bridge the gap between transformer architectures and earlier successes achieved in character-based language modeling, \citet{al2019character} applied extremely deep transformers and incorporated auxiliary loss functions to successfully train character-based transformer models that exceeded the performance of LSTM-driven character language models. Despite these improvements, a noticeable performance disparity compared to word-level large language models persisted. Similarly, GPT-2~\citep{radfordgpt2} found that when trained on large-scale corpora such as the 1 Billion Word Benchmark, language models operating at the byte level fell short of the results produced by those using word-level representations.

Although \cite{Choe2019BridgingTG} demonstrated that transformer-based byte-level LLMs can exceed the performance of similarly sized subword-level LLMs, this advantage comes at the cost of significantly greater computational demand and substantially longer training durations. In parallel, \cite{el2020characterbert} introduced a BERT variant (CharFormer) that generates word representations via convolutional operations over character embeddings, achieving gains within the medical domain, but their training process also requires considerably more computation. Building upon these ideas, \cite{clark2022canine} proposed CANINE, an encoder-only architecture with 150M parameters that works directly on character sequences. The model’s core is comprised of a deep transformer stack, analogous to the global module in our own design, and it incorporates both a local transformer and strided convolutional layers to reduce the resolution of the character inputs. This architecture surpasses its token-level encoder-only counterpart (mBERT) in downstream multilingual evaluation settings. ByT5~\citep{xue2022byt5} investigated encoder-decoder architectures at the byte level, deliberately avoiding patching mechanisms. Their model demonstrates greater noise robustness and achieves performance near that of tokenization-based models with just a quarter of the required data. However, in the absence of patching, their architecture must process computationally intensive attention across every byte, making resource usage substantially higher. The shift from subword tokens to raw byte modeling increases input sequence length, complicating the efficient scaling of byte-level architectures. More recently, advances with the Mamba Architecture~\citep{gu2023mamba}, which supports fixed-size memory states even for lengthy contexts, enabled \cite{wang2024mambabyte} to train a byte-level Mamba model without patching. At the 350M parameter level and under flops-matched settings, their approach surpasses byte-level transformer models in bits-per-byte across multiple datasets.

\textbf{Patching-based approaches:} Employing patching methods can significantly reduce the computational cost associated with byte-level LLMs, without necessarily compromising on performance. Several studies have showcased promising outcomes on relatively modest model sizes and limited training data. \cite{nawrot-etal-2022-hierarchical} introduce a static patching strategy utilizing downsampling and upsampling, culminating in the hourglass transformer, which demonstrates superior results compared to other byte-level models at the 150M parameter level. Building upon this, \cite{nawrot-etal-2023-efficient} advance these ideas by incorporating dynamic patching mechanisms, such as a boundary predictor trained in an end-to-end manner, a boundary predictor guided by specific tokenizers, and an entropy-based patching model reminiscent of {\textsc BLT}. Their experiments indicate that these refinements allow the approach to surpass standard transformers on language modeling benchmarks when using approximately 40M parameters trained on ~400M tokens. Further, \cite{lester2024training} explore the benefits of training on sequences compressed via arithmetic coding, achieving higher compression ratios than BPE. Through the equal-info windows technique, they report improved performance over byte-level baselines on language modeling tasks, though these gains do not extend to subword-based baselines.

Our research is primarily informed by, and bears the strongest connection to, MegaByte~\citep{yu2023megabyte}, a causal decoder-only language model that relies on a predetermined static patching approach. This method transforms raw bytes into patches by concatenating representations and subsequently utilizes a local model within the decoder to reconstruct bytes from these patches. The authors of MegaByte have shown that their approach can achieve performance on par with models based on tokenization when scaled to 1B parameters and evaluated on a dataset comprising approximately 400B bytes. Across all our experiments, we include ablations of MegaByte and observe that static patching is outperformed by contemporary compute-efficient tokenizer-based models in flop-matched evaluations. We further illustrate how {\textsc BLT} effectively narrows this performance divide. Independently, \cite{slagle2024spacebyte} arrived at a similar conclusion, and propose enhancing static patching by targeting whitespaces and analogous characters, as well as integrating a local encoder component. Their findings suggest gains over token-based transformers under compute-constrained conditions in specific settings, notably for Github and arXiv datasets at the 1B parameter level. We replicate experiments using their approach and demonstrate that advancing byte-level model architectures remains necessary in order to reliably achieve parity with leading token-based models such as Llama 3 at larger scales.

\section{Limitations and Future Work}

In this study, we adopt the optimal training step counts identified for Llama~3~\citep{dubey2024llama} to guide our architectural decisions. It is important to note, however, that these scaling laws were originally derived for BPE-level transformers, which may not yield ideal ratios between data and model parameters for {\textsc BLT}. The development of scaling laws tailored specifically to {\textsc BLT} remains an open direction for future work and could reveal even more advantageous scaling properties for our design. Furthermore, a significant portion of our experiments is conducted at a scale of up to 1B parameters; as we extend our models to 8B parameters and higher, the optimal architectural configuration may shift, potentially resulting in further performance gains at these larger scales.

Current transformer libraries and codebases prioritize efficiency for architectures based on tokenization. Although our study includes theoretical flop-equivalent experiments and leverages efficient solutions like FlexAttention for layers differing from the standard transformer structure, our implementations may still fall short of tokenizer-based models when considering actual wall-clock performance, indicating room for additional optimization.

Although {\textsc BLT} currently employs a separately trained entropy model for patching, an intriguing avenue for future research involves developing the patching model through end-to-end training. Section \ref{sec:distillation} outlines preliminary experiments, which provide promising evidence for the effectiveness of transforming tokenizer-based models such as Llama 3—which have been trained on over 10T tokens—into ``byte-ified'' models by initializing the global transformer with their pre-trained weights and keeping these weights fixed. Pursuing this research direction could reveal approaches that preserve the strengths of bytefying while potentially surpassing the performance of the original tokenizer-based models, all without the need for retraining from scratch.

\section{Conclusion}
\label{section:conclusion}

This work introduces the Byte Latent Transformer (\textbf{BLT}), an innovative model architecture that challenges the traditional reliance on predetermined tokenization schemes in large language models. By implementing a flexible, trainable mechanism for aggregating bytes into patches, {\textsc BLT} adaptively distributes computational effort in accordance with the complexity of the input, which yields marked gains in both computational efficiency and model robustness. Through comprehensive scaling experiments, we show that {\textsc BLT} achieves performance parity with established tokenization-dependent models such as Llama 3 at scales reaching 8B parameters and 4T bytes of data, and can attain up to a 50\% reduction in inference computational requirements with only slight diminutions in evaluation scores. Additionally, {\textsc BLT} enables a novel scaling strategy by facilitating concurrent increases in both patch size and model size, all while maintaining constant inference costs—a property particularly beneficial in compute-constrained real-world environments. By processing inputs directly at the byte level, {\textsc BLT} not only expands its resilience to noisy or atypical data but also enhances its capacity for fine-grained representation of sub-word components. Taken together, these findings establish {\textsc BLT} as a compelling, flexible, and robust alternative to conventional tokenization-based methods, offering an efficient and adaptable foundation for next-generation language models.

\end{document}